%%
%% Ergänzung zu Kapitel 6: Duale Registrierung im Studienbetrieb
%%

\begin{neu}
\section{Registrierungskanäle im Studienbetrieb}

Für die Durchführung der Studie stehen zwei administrative Zugangswege zur Verfügung. Direkt geworbene Interessierte registrieren sich mit ihrer E-Mail-Adresse und den für die spätere Auszahlung erforderlichen Angaben. Zusätzlich können über Prolific rekrutierte Personen ihre 24-stellige Prolific-ID als Teilnehmeridentifikator verwenden. Die Wahl des Registrierungskanals verändert die experimentelle Durchführung nicht: In beiden Fällen gelten dieselben inhaltlichen Einschlusskriterien, die produktive App verwendet denselben App-Token-basierten Studienzugang und es sind dieselben zwanzig Selbstberichte mit zehn Cue-Matching- und zehn Cue-Labeling-Situationen vorgesehen.

Der direkte Teilnahmeweg umfasst nach dem Absenden des Webformulars weiterhin ein E-Mail-Double-Opt-In. Erst nach Aufruf des Bestätigungslinks kann die Registrierung für die App-Aktivierung verwendet werden. Bei einer Prolific-Registrierung werden Name und Bankverbindung im CueLens-Formular nicht erhoben und es wird kein zusätzlicher E-Mail-Bestätigungsprozess durchgeführt. Die Aktivierung erfolgt anschließend mit derselben App-Oberfläche, die sowohl eine E-Mail-Adresse als auch eine Prolific-ID akzeptiert. Der technische App-Token-Handshake ist für beide Wege gleich aufgebaut und führt nach erfolgreicher Bestätigung in denselben produktiven Studienzustand.

Die kanalabhängige Erhebung administrativer Angaben reduziert die Teilnahmebelastung für über Prolific rekrutierte Personen und vermeidet eine doppelte Speicherung von Daten, die für diesen Zugangsweg nicht benötigt werden. Die Prolific-ID verbleibt dabei in der Registrierungsdatenbank. Sie wird nicht zusammen mit den Craving-Werten an den wissenschaftlichen Selbstberichtsendpunkt übertragen. Damit bleibt der primäre Studiendatensatz unabhängig davon strukturiert, über welchen Rekrutierungskanal eine Person in die Studie gelangt ist.

Für die Dokumentation der Rekrutierung ist der Kanal dennoch relevant. Die Zahl der Anmeldungen und vollständigen Teilnahmen sollte getrennt nach direkter Rekrutierung und Prolific berichtet werden, weil sich die beiden Zugangswege hinsichtlich Selbstselektion, technischer Vorerfahrung und Anreizstruktur unterscheiden können. Eine solche deskriptive Darstellung dient der Einordnung der Stichprobe; sie verändert nicht den vorab definierten Within-Subject-Vergleich der Craving-Werte. Die externe Prüfung einer Prolific-ID auf Zugehörigkeit zur CueLens-Studie und die abweichende administrative Behandlung des Studienabschlusses werden in den nachfolgenden technischen Ergänzungen gesondert beschrieben.
\end{neu}
